% !TEX program = xelatex
\documentclass[10pt,a4paper]{article}
\usepackage[T1]{fontenc}
\usepackage[french]{babel}
\usepackage[left=0.6cm,right=0.6cm,top=1cm,bottom=1cm]{geometry}
\usepackage{amsmath}
\usepackage{amssymb}
\usepackage{array}
\usepackage{graphicx}
\usepackage{xcolor}
\usepackage{titlesec}
\usepackage{enumitem}
\usepackage{multicol}
\usepackage{pifont}
\usepackage{frcursive}
% Correctif du bug de crénage de frcursive sous XeLaTeX (repris de CLT2627.sty) :
% sans lui, certaines paires de lettres consécutives dans \cursive explosent
% (une lettre projetée loin à droite), ce qui coupe des mots comme "Représenter".
\makeatletter
\ifdefined\XeTeXcharclass
\XeTeXcharclass `a 201 \XeTeXcharclass `b 201 \XeTeXcharclass `c 201
\XeTeXcharclass `d 201 \XeTeXcharclass `e 201 \XeTeXcharclass `f 201
\XeTeXcharclass `g 201 \XeTeXcharclass `h 201 \XeTeXcharclass `i 201
\XeTeXcharclass `j 201 \XeTeXcharclass `k 201 \XeTeXcharclass `l 201
\XeTeXcharclass `m 201 \XeTeXcharclass `n 201 \XeTeXcharclass `o 201
\XeTeXcharclass `p 201 \XeTeXcharclass `q 201 \XeTeXcharclass `r 201
\XeTeXcharclass `s 201 \XeTeXcharclass `t 201 \XeTeXcharclass `u 201
\XeTeXcharclass `v 201 \XeTeXcharclass `w 201 \XeTeXcharclass `x 201
\XeTeXcharclass `y 201 \XeTeXcharclass `z 201
\XeTeXcharclass `A 201 \XeTeXcharclass `B 201 \XeTeXcharclass `C 201
\XeTeXcharclass `D 201 \XeTeXcharclass `E 201 \XeTeXcharclass `F 201
\XeTeXcharclass `G 201 \XeTeXcharclass `H 201 \XeTeXcharclass `I 201
\XeTeXcharclass `J 201 \XeTeXcharclass `K 201 \XeTeXcharclass `L 201
\XeTeXcharclass `M 201 \XeTeXcharclass `N 201 \XeTeXcharclass `O 201
\XeTeXcharclass `P 201 \XeTeXcharclass `Q 201 \XeTeXcharclass `R 201
\XeTeXcharclass `S 201 \XeTeXcharclass `T 201 \XeTeXcharclass `U 201
\XeTeXcharclass `V 201 \XeTeXcharclass `W 201 \XeTeXcharclass `X 201
\XeTeXcharclass `Y 201 \XeTeXcharclass `Z 201
\XeTeXcharclass `\é 201 \XeTeXcharclass `\è 201 \XeTeXcharclass `\ê 201
\XeTeXcharclass `\ë 201 \XeTeXcharclass `\É 201 \XeTeXcharclass `\È 201
\XeTeXcharclass `\Ê 201 \XeTeXcharclass `\Ë 201 \XeTeXcharclass `\à 201
\XeTeXcharclass `\â 201 \XeTeXcharclass `\ä 201 \XeTeXcharclass `\À 201
\XeTeXcharclass `\Â 201 \XeTeXcharclass `\Ä 201 \XeTeXcharclass `\ô 201
\XeTeXcharclass `\ö 201 \XeTeXcharclass `\Ô 201 \XeTeXcharclass `\Ö 201
\XeTeXcharclass `\î 201 \XeTeXcharclass `\ï 201 \XeTeXcharclass `\Î 201
\XeTeXcharclass `\Ï 201 \XeTeXcharclass `\ù 201 \XeTeXcharclass `\û 201
\XeTeXcharclass `\ü 201 \XeTeXcharclass `\Ù 201 \XeTeXcharclass `\Û 201
\XeTeXcharclass `\Ü 201 \XeTeXcharclass `\ç 201 \XeTeXcharclass `\Ç 201
\XeTeXcharclass `\œ 201 \XeTeXcharclass `\Œ 201 \XeTeXcharclass `\æ 201
\XeTeXcharclass `\Æ 201
\let\cursive@orig\cursive
\renewcommand{\cursive}{%
  \cursive@orig
  \XeTeXinterchartoks 201 201 = {\kern0pt }%
}
\fi
\makeatother
\usepackage{tikz}
\usetikzlibrary{calc,arrows.meta,shapes.misc}
\tikzset{ptcross/.style={cross out,draw,thick,minimum size=3pt,inner sep=0pt}}
\usepackage{tcolorbox}
\tcbuselibrary{skins,breakable}
\usepackage[hidelinks]{hyperref}
\usepackage{fontawesome5}
\usepackage{qrcode}
% Rendre un \qrcode* cliquable : boîte vide de la taille du QR, superposée
% via \href (même astuce que dans le Cours).
\newcommand{\ClickableQR}[2]{\makebox[0pt][l]{\href{#2}{\makebox[#1][c]{\rule{0pt}{#1}}}}}
\definecolor{calcblue}{HTML}{1C7ED6}
\definecolor{diffgreen}{HTML}{40C057}
\definecolor{diffblue}{HTML}{228BE6}
\definecolor{diffred}{HTML}{FA5252}
\definecolor{skyblue}{HTML}{339AF0}
\definecolor{objorange}{HTML}{F76707}
\newcommand{\starfull}[1]{{\color{#1}\ding{72}}}
\newcommand{\starempty}{{\color{gray!45}\ding{73}}}
\newcommand{\Vj}[1]{\overrightarrow{\text{#1}}}
\newcommand{\VBlank}{\rule[0.15ex]{0.55cm}{0.35pt}}
\newcommand{\Pt}[1]{\textsf{\footnotesize #1}}
\renewcommand{\familydefault}{\sfdefault}
\DeclareSymbolFont{sfnumbers}{OT1}{cmss}{m}{n}
\SetSymbolFont{sfnumbers}{bold}{OT1}{cmss}{bx}{n}
\DeclareMathSymbol{0}{\mathalpha}{sfnumbers}{`0}
\DeclareMathSymbol{1}{\mathalpha}{sfnumbers}{`1}
\DeclareMathSymbol{2}{\mathalpha}{sfnumbers}{`2}
\DeclareMathSymbol{3}{\mathalpha}{sfnumbers}{`3}
\DeclareMathSymbol{4}{\mathalpha}{sfnumbers}{`4}
\DeclareMathSymbol{5}{\mathalpha}{sfnumbers}{`5}
\DeclareMathSymbol{6}{\mathalpha}{sfnumbers}{`6}
\DeclareMathSymbol{7}{\mathalpha}{sfnumbers}{`7}
\DeclareMathSymbol{8}{\mathalpha}{sfnumbers}{`8}
\DeclareMathSymbol{9}{\mathalpha}{sfnumbers}{`9}
\setlength{\parindent}{0pt}
\graphicspath{{./}}
\newcommand*{\textcircledd}[1]{\begin{tikzpicture}[baseline=(letter.base)]\node[draw,circle,inner sep=1pt] (letter) {#1};\end{tikzpicture}}
\newcommand{\numeroCerne}[1]{\protect\textcircledd{\bfseries\scriptsize #1}}
\setlist[enumerate,1]{label=\numeroCerne{\arabic*}, leftmargin=*, labelsep=6pt, itemsep=3pt, topsep=2pt, partopsep=0pt, parsep=0pt}

\NewDocumentEnvironment{Exo}{O{breakable} m m m}{%
  \begin{tcolorbox}[
    enhanced,
    #1,
    colback=white,
    colframe=white,
    boxrule=0pt,
    arc=4pt, outer arc=4pt, sharp corners=south,
    left=8pt, right=8pt, top=5pt, bottom=3pt,
    before skip=4pt, after skip=4pt,
  ]
  \noindent\tcbox[on line, colback=skyblue, colframe=skyblue, boxsep=1pt, left=4pt, right=4pt, top=1pt, bottom=1pt, arc=2pt]{\textcolor{white}{\bfseries Exercice~#2}}~#3\hfill\mbox{\color{skyblue}\normalsize\cursive\hyphenchar\font=-1 #4}\par\vspace{3pt}%
}{%
  \end{tcolorbox}
}

\NewDocumentEnvironment{ExoCont}{O{unbreakable}}{%
  \begin{tcolorbox}[
    enhanced,
    #1,
    colback=white,
    colframe=white,
    boxrule=0pt,
    arc=4pt, outer arc=4pt, sharp corners=south,
    left=8pt, right=8pt, top=5pt, bottom=3pt,
    before skip=4pt, after skip=4pt,
  ]
}{%
  \end{tcolorbox}
}

\newcommand{\Objectif}[2]{%
\vspace{4pt}
\noindent{\large\color{objorange}\bfseries Objectif #1~:}\\{\large\color{objorange}#2}\par
\vspace{2pt}
}

\NewDocumentCommand{\FeuilletTitre}{m}{%
  \begin{center}
  \begin{tcolorbox}[
    enhanced,
    width=0.58\linewidth,
    colback=white,
    colframe=black,
    boxrule=0.8pt,
    arc=8pt,
    sharp corners=downhill,
    left=10pt, right=10pt, top=8pt, bottom=6pt,
    halign=center,
    overlay={
      \node[fill=white, anchor=west, font=\normalsize, inner xsep=4pt]
        at ([xshift=14pt]frame.north west) {Feuille d'exercices n\textdegree{}3\,\textsuperscript{3/4}};
    },
  ]
  {\huge #1}
  \end{tcolorbox}
  \end{center}
}

% Trous "texte à trous" : vecteur dont le nom est à compléter, fraction à
% compléter, réel à compléter.
\newcommand{\Vd}{\overrightarrow{\ \dots\dots\ }}
% Variante avec un espace (strut invisible) entre les pointillés et la
% flèche, pour pouvoir écrire le nom du vecteur sous la flèche.
\newcommand{\VdE}{\overrightarrow{\rule{0pt}{2.6ex}\ \dots\dots\ }}
\newcommand{\VdEl}[1]{\overrightarrow{\rule{0pt}{2.6ex}\ \dots\,\text{#1}\ }}
\newcommand{\Fd}{\dfrac{\ \dots\ }{\ \dots\ }}
\newcommand{\Nd}{\dots\,}
% Ligne pointillée pour rédiger une justification.
\newcommand{\LigneJ}{\noindent\phantom{.}\dotfill\par\vspace{13pt}}
\setlength{\abovedisplayskip}{3pt}\setlength{\belowdisplayskip}{3pt}\setlength{\abovedisplayshortskip}{3pt}\setlength{\belowdisplayshortskip}{3pt}

\pagestyle{empty}
\begin{document}

\FeuilletTitre{Vecteurs et d\'emonstrations}

\Objectif{5}{Produit d'un vecteur par un nombre r\'eel et colin\'earit\'e de deux vecteurs}

% ---------------------------------------------------------------------
% Ex 1 : démonstration d'alignement, guidée (ancienne Feuille 6 - Ex 2)
% ---------------------------------------------------------------------
\begin{Exo}[breakable]{1}{\starfull{diffblue}\starfull{diffblue}\starempty}{Raisonner}
\noindent\begin{minipage}[t]{0.62\linewidth}\vspace{0pt}
Soit $ABC$ un triangle. On d\'efinit les points $M$ et $N$ par les \'egalit\'es vectorielles suivantes~: \quad $\Vj{BM}=2\,\Vj{AC}$ \quad et \quad $\Vj{AN}=\dfrac12\,\Vj{BC}$.\\[6pt]
\numeroCerne{1}~Placer les points $M$ et $N$ sur la figure ci-contre.\\[6pt]
\textbf{Observation}~: \dotfill\\[4pt]
\numeroCerne{2}~On va chercher \`a montrer que les points $C$, $M$ et $N$ sont align\'es. Pour cela, on exprime les vecteurs $\Vj{CM}$ et $\Vj{CN}$ en fonction des vecteurs $\Vj{AC}$ et $\Vj{CB}$.
\end{minipage}\hfill
\begin{minipage}[t]{0.36\linewidth}\vspace{0pt}
\centering
\begin{tikzpicture}[scale=0.55]
\draw[thin,gray!35] (-2,-1) grid (6,5);
\coordinate (A) at (0,0); \coordinate (B) at (3,0); \coordinate (C) at (1,2);
\draw[thick] (A)--(B)--(C)--cycle;
\node[ptcross] at (A) {}; \node[below left] at (A) {\Pt{A}};
\node[ptcross] at (B) {}; \node[below right] at (B) {\Pt{B}};
\node[ptcross] at (C) {}; \node[above] at (C) {\Pt{C}};
\end{tikzpicture}
\end{minipage}

\medskip

\noindent\begin{minipage}[t]{0.5\linewidth}\vspace{0pt}
\textbf{D'une part,}
\[\begin{aligned}
\Vj{CM} &= \VdEl{B} + \VdE \quad\text{(relation de \dots\dots\dots\dots)}\\
        &= \Vj{CB} + 2\,\Vj{AC}
\end{aligned}\]
\end{minipage}\hfill
\begin{minipage}[t]{0.48\linewidth}\vspace{0pt}
\textbf{D'autre part,}
\[\begin{aligned}
\Vj{CN} &= \VdEl{A} + \VdE \quad\text{(relation de \dots\dots\dots\dots)}\\
        &= -\VdE + \dfrac12\,\VdE \\
        &= -\dfrac12\,\VdE - \VdE
\end{aligned}\]
\end{minipage}

\medskip

\numeroCerne{3}~On remarque que~: \fbox{\rule[-1.2ex]{0pt}{4ex}$\Vj{CM} = \Nd\Nd\ \Vj{CN}$}\quad
donc les vecteurs $\VdE$ et $\VdE$ sont \dotfill

\medskip

\fbox{\rule[-0.9ex]{0pt}{3.2ex}\textbf{Ainsi}~: les points \Nd\Nd, \Nd\Nd\ et \Nd\Nd\ sont \dots\dots\dots\dots\dots\dots\dots\dots\dots\dots\dots\ .}
\end{Exo}

% ---------------------------------------------------------------------
% Ex 2 : démonstration de parallélisme, guidée (ancienne Feuille 6 - Ex 1)
% NB : avec AM = 1/4 AB et BN = 3/4 BC, on obtient MN = 3/4 AC, donc c'est
% (MN) // (AC) que l'on démontre (l'ancienne feuille annonçait (BC) par
% erreur, mais le calcul à trous menait bien à AC).
% ---------------------------------------------------------------------
\begin{Exo}[breakable]{2}{\starfull{diffblue}\starfull{diffblue}\starempty}{Raisonner}
\noindent\begin{minipage}[t]{0.62\linewidth}\vspace{0pt}
Soit $ABC$ un triangle. On d\'efinit les points $M$ et $N$ par les \'egalit\'es vectorielles suivantes~: \quad $\Vj{AM}=\dfrac14\,\Vj{AB}$ \quad et \quad $\Vj{BN}=\dfrac34\,\Vj{BC}$.\\[6pt]
\numeroCerne{1}~Placer les points $M$ et $N$ sur la figure ci-contre.\\[6pt]
\textbf{Observation}~: \dotfill\\[4pt]
\numeroCerne{2}~On va chercher \`a montrer que les droites $(MN)$ et $(AC)$ sont parall\`eles. Pour cela, on va montrer que les vecteurs $\Vj{MN}$ et $\Vj{AC}$ sont colin\'eaires, c'est-\`a-dire chercher un r\'eel $k$ tel que $\Vj{MN}=k\,\Vj{AC}$.
\end{minipage}\hfill
\begin{minipage}[t]{0.36\linewidth}\vspace{0pt}
\centering
\begin{tikzpicture}[scale=0.55]
\draw[thin,gray!35] (-1,-1) grid (7,5);
\coordinate (A) at (0,0); \coordinate (B) at (4,0); \coordinate (C) at (6,4);
\draw[thick] (A)--(B)--(C)--cycle;
\node[ptcross] at (A) {}; \node[below left] at (A) {\Pt{A}};
\node[ptcross] at (B) {}; \node[below right] at (B) {\Pt{B}};
\node[ptcross] at (C) {}; \node[above] at (C) {\Pt{C}};
\end{tikzpicture}
\end{minipage}

\medskip

\noindent\begin{minipage}[t]{0.5\linewidth}\vspace{0pt}
\textbf{Tout d'abord}, exprimons $\Vj{MB}$ en fonction de $\Vj{AB}$~:
\[\begin{aligned}
\Vj{MB} &= \overrightarrow{\rule{0pt}{2.6ex}\ \text{M}\,\dots\ } + \Vj{AB} \\
        &= \Fd\,\Vj{AB} + \Vj{AB} \\
        &= \Fd\,\Vj{AB}
\end{aligned}\]
\end{minipage}\hfill
\begin{minipage}[t]{0.46\linewidth}\vspace{0pt}
\textbf{Justifications}~:\\[6pt]
\LigneJ
\LigneJ
\end{minipage}

\medskip

\noindent\begin{minipage}[t]{0.5\linewidth}\vspace{0pt}
\textbf{Ensuite},
\[\begin{aligned}
\Vj{MN} &= \Vj{MB} + \VdE \\
        &= \Nd\,\Vj{AB} + \Fd\,\VdE \\
        &= \Fd\left(\VdE + \VdE\right) \\
        &= \Fd\,\VdE
\end{aligned}\]
\end{minipage}\hfill
\begin{minipage}[t]{0.46\linewidth}\vspace{0pt}
\textbf{Justifications}~:\\[6pt]
\LigneJ
\LigneJ
\LigneJ
\end{minipage}

\medskip

\textbf{Conclusion}~: les vecteurs $\VdE$ et $\VdE$ sont \dotfill

\medskip

\fbox{\rule[-0.9ex]{0pt}{3.2ex}\textbf{Ainsi}~: les droites $(\Nd\Nd)$ et $(\Nd\Nd)$ sont \dots\dots\dots\dots\dots\dots\dots\dots\dots\dots\dots\ .}
\end{Exo}

\newpage

% ---------------------------------------------------------------------
% Page 2 : démonstrations à rédiger entièrement.
% ---------------------------------------------------------------------
\begin{tcolorbox}[enhanced, colback=objorange!6, colframe=objorange, boxrule=0.8pt, arc=4pt,
  left=8pt, right=8pt, top=5pt, bottom=5pt, before skip=4pt, after skip=8pt,
  title={\bfseries M\'ethode -- d\'emontrer avec des vecteurs}, fonttitle=\bfseries, coltitle=white, colbacktitle=objorange]
\begin{itemize}[label=$\star$, leftmargin=*, itemsep=2pt]
\item Les droites $(AB)$ et $(CD)$ sont \textbf{parall\`eles} $\Longleftrightarrow$ les vecteurs $\Vj{AB}$ et $\Vj{CD}$ sont colin\'eaires.
\item Les points $A$, $B$ et $C$ sont \textbf{align\'es} $\Longleftrightarrow$ les vecteurs $\Vj{AB}$ et $\Vj{AC}$ sont colin\'eaires.
\item $I$ est le \textbf{milieu} de $[AB]$ $\Longleftrightarrow$ $\Vj{AI}=\Vj{IB}$ $\Longleftrightarrow$ $\Vj{IA}+\Vj{IB}=\vec0$ $\Longleftrightarrow$ $\Vj{AI}=\dfrac12\,\Vj{AB}$.
\end{itemize}
Dans chaque cas, on exprime les vecteurs concern\'es en fonction de \textbf{deux m\^emes vecteurs} (relation de Chasles, \'egalit\'es de l'\'enonc\'e), puis on cherche le r\'eel $k$ qui les relie.
\end{tcolorbox}

\begin{Exo}[unbreakable]{3}{\starfull{diffblue}\starfull{diffblue}\starempty}{Raisonner}
\noindent\begin{minipage}[t]{0.6\linewidth}\vspace{0pt}
Soit $ABC$ un triangle. On d\'efinit les points $M$ et $N$ par~:
\[\Vj{AM}=-\dfrac12\,\Vj{AB}\qquad\text{et}\qquad\Vj{AN}=-\dfrac12\,\Vj{AC}.\]
\begin{enumerate}
\item Compl\'eter la figure ci-contre.
\item Exprimer le vecteur $\Vj{MN}$ en fonction du vecteur $\Vj{BC}$.
\item En d\'eduire que les droites $(MN)$ et $(BC)$ sont parall\`eles.
\end{enumerate}
\end{minipage}\hfill
\begin{minipage}[t]{0.38\linewidth}\vspace{0pt}
\centering
\begin{tikzpicture}[scale=0.55]
\draw[thin,gray!35] (0,0) grid (9,5);
\coordinate (A) at (4,2); \coordinate (B) at (8,4); \coordinate (C) at (2,4);
\draw[thick] (A)--(B)--(C)--cycle;
\node[ptcross] at (A) {}; \node[below left] at (A) {\Pt{A}};
\node[ptcross] at (B) {}; \node[below right] at (B) {\Pt{B}};
\node[ptcross] at (C) {}; \node[above left] at (C) {\Pt{C}};
\end{tikzpicture}
\end{minipage}
\end{Exo}

\begin{Exo}[unbreakable]{4}{\starfull{diffred}\starfull{diffred}\starfull{diffred}}{Raisonner}
\noindent\begin{minipage}[t]{0.78\linewidth}\vspace{0pt}
Soit $ABCD$ un parall\'elogramme. On d\'efinit les points $E$ et $F$ par~:
\[\Vj{BE}=\dfrac12\,\Vj{BC}\qquad\text{et}\qquad\Vj{DF}=2\,\Vj{DC}.\]
\begin{enumerate}
\item R\'ealiser la figure sur GeoGebra~: \faLaptop~\textbf{Code Capytale}~: \href{https://capytale2.ac-paris.fr/web/c/414a-11688973/edn}{{\ttfamily\bfseries 414a-11688973}}\\[2pt]
Que peut-on conjecturer pour les points $A$, $E$ et $F$~?
\item Exprimer le vecteur $\Vj{AE}$ en fonction des vecteurs $\Vj{AB}$ et $\Vj{AD}$.
\item Exprimer le vecteur $\Vj{AF}$ en fonction des vecteurs $\Vj{AB}$ et $\Vj{AD}$.
\item D\'emontrer que les points $A$, $E$ et $F$ sont align\'es.
\end{enumerate}
\end{minipage}\hfill
\begin{minipage}[t]{0.2\linewidth}\vspace{0pt}
\centering
\ClickableQR{2cm}{https://capytale2.ac-paris.fr/web/c/414a-11688973/edn}\textcolor{black}{\qrcode*[height=2cm, level=M]{https://capytale2.ac-paris.fr/web/c/414a-11688973/edn}}
\end{minipage}
\end{Exo}

\Objectif{7}{Caract\'erisation vectorielle du milieu d'un segment}

\begin{Exo}[unbreakable]{5}{\starfull{diffred}\starfull{diffred}\starfull{diffred}}{Raisonner}
Soit $ABCD$ un parall\'elogramme. On note $E$ le sym\'etrique de $A$ par rapport \`a $B$, et $F$ le sym\'etrique de $A$ par rapport \`a $D$.
\begin{enumerate}
\item Faire une figure.
\item Exprimer le vecteur $\Vj{CE}$ en fonction des vecteurs $\Vj{AB}$ et $\Vj{AD}$.
\item Exprimer le vecteur $\Vj{CF}$ en fonction des vecteurs $\Vj{AB}$ et $\Vj{AD}$.
\item En d\'eduire que $C$ est le milieu du segment $[EF]$.
\end{enumerate}
\end{Exo}

\end{document}
